\chapter{State of the Art}
\label{ch:sota}

Link prediction methods can be grouped into three families, followed by the specific line of work on
attention in multiplex networks that this report builds on. Recent work gives a broad
overview of the area~\cite{zangari2024,heuristickdd2024}.

\section{Classical Similarity Methods}
The earliest and still widely used approach scores a candidate pair from the structure of its shared
neighbourhood. \emph{Common Neighbours} counts shared neighbours; \emph{Jaccard} normalises that
count; \emph{Adamic--Adar}~\cite{adamic2003} down-weights popular shared neighbours; and
\emph{Preferential Attachment} multiplies the two node degrees. These heuristics are fast, require no
training, and remain strong baselines, and they continue to be refined~\cite{heuristickdd2024}. Their
defining limitation is structural: when a pair shares \emph{no} neighbours -- exactly the weak-tie
case -- every neighbourhood score is zero and the prediction is no better than chance.

\section{Embedding-based Methods}
A second family learns a low-dimensional vector for each node such that nodes with similar roles lie
close together. \emph{Node2Vec}~\cite{node2vec} generates these vectors from biased random walks and
predicts a link from the similarity of the two node vectors. Because the vectors capture global
role information rather than only immediate overlap, such methods can score pairs that share no
neighbours, and they achieve strong ranking performance. They do not, however, use node attributes
and, in their basic form, treat each layer separately.

\section{Graph Neural Networks and Attention}
Graph neural networks learn node representations by repeatedly aggregating information from
neighbours. The \emph{Graph Convolutional Network} (GCN)~\cite{gcn} averages neighbour features with
fixed, degree-based weights. The \emph{Graph Attention Network} (GAT)~\cite{gat} replaces this fixed
average with a learned \emph{attention} mechanism, so that each neighbour receives an importance
weight conditioned on its features; \emph{GATv2}~\cite{gatv2} later showed that a small
reordering of GAT's operations yields a strictly more expressive, \emph{dynamic} attention.
Attention has since become central to graph learning, both in dedicated
surveys~\cite{attnsurvey2023} and in the broader move toward graph
transformers~\cite{gtsurvey2024}. Attention matters here for two reasons: it improves
accuracy by focusing on informative neighbours, and it produces weights that can be read back as an
\emph{explanation} of the prediction.

\section{Attention in Multiplex Networks}
Extending attention to multilayer data is an active direction. The \emph{Heterogeneous Attention
Network} (HAN)~\cite{han2019} introduced two levels of attention -- within a relation and across
relations -- and is the conceptual starting point for the cross-layer attention used in this work.
More recent studies apply attention and message passing directly to multiplex link
prediction~\cite{multiplexsage2024,gao2023}, including cross-layer attention mechanisms closely
related to the one used here.

\section{Research Gap}
Across these families, two gaps persist. First, the methods that are strongest on ordinary links --
the neighbourhood heuristics -- are by construction random on weak-tie links, while the methods that
can handle weak ties rarely target them explicitly. Second, cross-layer information is often either
ignored (single-layer models) or fused in a fixed, non-interpretable way. This work addresses both:
it injects a weak-tie signal into an attention encoder and adds an interpretable cross-layer
attention that learns how much each layer contributes.
